% specimen.tex -- one page that proves the design flags work (XeLaTeX).
%
% Built for every package and ground by design/specimen/build_specimens.sh:
%   xelatex "\def\tpcspecpkg{3}\def\tpcspecground{dark}% specimen.tex -- one page that proves the design flags work (XeLaTeX).
%
% Built for every package and ground by design/specimen/build_specimens.sh:
%   xelatex "\def\tpcspecpkg{3}\def\tpcspecground{dark}% specimen.tex -- one page that proves the design flags work (XeLaTeX).
%
% Built for every package and ground by design/specimen/build_specimens.sh:
%   xelatex "\def\tpcspecpkg{3}\def\tpcspecground{dark}% specimen.tex -- one page that proves the design flags work (XeLaTeX).
%
% Built for every package and ground by design/specimen/build_specimens.sh:
%   xelatex "\def\tpcspecpkg{3}\def\tpcspecground{dark}\input{specimen.tex}"
% Every colour, face and length below is a token name from design/README.md.
% This file contains no package-specific value; design/lint_content.py checks that.
%
% Text and physics are taken from the lane specimen
% docs/plan/research/design-packages-A/specimens/body-de.tex (checked there by hand):
%   rho = (1 + r.sigma)/2, E = (1 + n.sigma)/2, |n| = 1  =>  tr(rho E) = (1 + r.n)/2.
% Figure geometry: r = 0.75 at 50 degrees, n = e_x, so r.n = 0.75 cos 50 = 0.482 and
% p = (1 + 0.482)/2 = 0.741.  Ladder: particle in a box, E_n = n^2 E_1
% (latex/lectures/de/vl11.tex, latex/figures/fig_vl11_kasten.tex), n = 1..6.
\providecommand\tpcspecpkg{1}
\providecommand\tpcspecground{light}
\providecommand\tpcspecmath{stix}
\providecommand\tpcspecfonts{open}
\documentclass[11pt]{article}
\usepackage[ngerman]{babel}
\usepackage[autostyle]{csquotes}
\usepackage{tikz}
\usetikzlibrary{arrows.meta}
\usepackage[pkg=\tpcspecpkg, ground=\tpcspecground, math=\tpcspecmath,
            fonts=\tpcspecfonts]{tpc-design}
\usepackage[most]{tcolorbox}
\usepackage[protrusion=true, expansion=false]{microtype}
\usepackage[hidelinks]{hyperref}
\tpcapplypage
\pagestyle{empty}
\emergencystretch=1em
\setlength{\parindent}{\tpcspacem}
\setlength{\parskip}{0pt}

% ---- semantic devices, built from tokens only ------------------------------
\newcommand\speckicker[1]{{\tpcfontlabel\color{tpc-ink-muted}#1}\par\vspace{\tpcspacexs}}
\newcommand\spechead[2]{%
  \noindent\llap{{\tpcfontnumeral\color{tpc-accent}#1}\hspace{\tpcspaces}}%
  {\tpcfonthi\color{tpc-accent}#2}\par\vspace{\tpcspaces}}
\newcommand\specsection[1]{\par\vspace{\tpcspaces}\noindent
  {\tpcfonthii\color{tpc-accent}#1}\par\vspace{\tpcspacexs}\noindent\ignorespaces}
\newcommand\specnote[1]{\marginpar{\raggedright\tpcfontlabel\color{tpc-ink-muted}#1}}
\newcommand\speclink[1]{{\color{tpc-link}\underline{#1}}}
% Merksatz: highlight tint, box rule in the accent, run-in label.
\newtcolorbox{specmerksatz}[1]{enhanced, frame hidden, boxrule=0pt, arc=0pt,
  colback=tpc-highlight, coltext=tpc-ink, left=\tpcspaces, right=\tpcspaces,
  top=\tpcspaces, bottom=\tpcspaces, before skip=\tpcspaces, after skip=\tpcspaces,
  borderline west={\tpcstrokebox}{0pt}{tpc-accent},
  boxsep=0pt,
  before upper={\begingroup\tpcfonthiii\color{tpc-accent}#1\endgroup\quad}}
% Depth marking, the same device in every package: body size, no tint, no small
% type; a rule in the outer margin over the block's length and a run-in label.
% Only the rule's colour and weight and the label's face come from the package.
% #1 label, #2 title, #3 pointer to the Fundament (set beside the block).
\newtcolorbox{specvertiefung}[3]{enhanced, frame hidden, boxrule=0pt, arc=0pt,
  colback=tpc-ground, coltext=tpc-ink, boxsep=0pt, left=0pt, right=0pt, top=0pt, bottom=0pt,
  before skip=\tpcspaces, after skip=\tpcspaces,
  overlay={%
    \fill[tpc-depth-v] ([xshift=\tpcspaces]frame.north east)
      rectangle ([xshift=\dimexpr\tpcspaces+\tpcstrokedepthv\relax]frame.south east);
    % Fundament pointer: double hairline in the depth-f role, then the label
    \coordinate (F) at ([xshift=\tpcmargingap, yshift=-3pt]frame.north east);
    \draw[tpc-depth-f, line width=\tpcstrokedepthf] (F) -- ++(\tpcmarginwidth,0)
      ([yshift=-2pt]F) -- ++(\tpcmarginwidth,0);
    \node[anchor=north west, inner sep=0pt, text width=\tpcmarginwidth, align=flush left,
      text=tpc-depth-f, font=\tpcfontlabel] at ([yshift=-5pt]F) {#3};},
  before upper={\begingroup\tpcfontdepth\color{tpc-depth-v}Vertiefung #1\enspace #2\par\endgroup}}
\newcommand\vr{\symbfit{r}}
\newcommand\vn{\symbfit{n}}
\newcommand\vsigma{\symbfit{\sigma}}
\ifdefined\symbfit\else
  \renewcommand\vr{\boldsymbol{r}}\renewcommand\vn{\boldsymbol{n}}
  \renewcommand\vsigma{\boldsymbol{\sigma}}
\fi
\newcommand\one{\mathbb{1}}

\renewcommand\theequation{6.\arabic{equation}}
\begin{document}
\speckicker{Zoomstufe 1 \textendash\ ein einzelnes Quantensystem}
\spechead{6}{Die Bornsche Regel}

\noindent Eine physikalische Theorie liefert Daten für drei Dinge: Präparationen, Dynamik und
Messungen. Eine Präparation nennen wir \enquote{Zustand}; im Quantenfall beschreiben wir sie
durch einen Dichteoperator $\rho\geq 0$ mit $\operatorname{tr}\rho=1$.\specnote{Für ein Qubit ist
jeder Dichteoperator von der Form $\rho=\tfrac12(\one+\vr\cdot\vsigma)$ mit
$\lvert\vr\rvert\leq 1$. Reine Zustände liegen auf dem Rand.}
Ein Detektor, der \enquote{klick} macht oder schweigt, wird durch einen Effekt $E$ mit
$0\leq E\leq\one$ beschrieben. Die Theorie sagt nur Wahrscheinlichkeiten voraus, und zwar genau
eine Zahl für jedes Paar aus Zustand und Effekt (siehe \speclink{Definition D-6.1}).

\begin{specmerksatz}{Postulat P4}
Wird ein System im Zustand $\rho$ präpariert und mit dem Effekt $E$ getestet, so spricht der
Detektor mit der Wahrscheinlichkeit $p(E\,\vert\,\rho)=\operatorname{tr}(\rho E)$ an.
\end{specmerksatz}

\noindent Für ein Qubit lässt sich die Regel zeichnen (Abb.~6.1a). Mit
$\rho=\tfrac12(\one+\vr\cdot\vsigma)$ und dem Projektor $E=\tfrac12(\one+\vn\cdot\vsigma)$,
$\lvert\vn\rvert=1$, folgt
\begin{equation}
  \operatorname{tr}(\rho E)=\tfrac12\bigl(1+\vr\cdot\vn\bigr),
  \qquad
  \rho\;\mapsto\;U_t\,\rho\,U_t^{*},\quad U_t=\mathrm{e}^{-\mathrm{i}Ht/\hbar}.
\end{equation}
Die Wahrscheinlichkeit ist also die Projektion des Bloch-Vektors auf die Richtung des Detektors,
verschoben und halbiert. Die Zeitentwicklung dreht den Vektor; sie ändert $\lvert\vr\rvert$ nicht.

% ---- figure: plate (a, b) and stage (c) -------------------------------------
\par\vspace{\tpcspaces}
\noindent\makebox[\linewidth][l]{%
\tpcplate{%
\begin{tikzpicture}[x=1cm, y=1cm, baseline=(current bounding box.south)]
  \useasboundingbox (-1.62,-1.78) rectangle (5.3,1.78);
  % (a) Bloch disc: state, effect, dynamics, bath
  \begin{scope}[scale=1.28]
    \fill[tpc bath, even odd rule] (0,0) circle (1.17) (0,0) circle (1.0);
    \fill[tpc bath hatch, even odd rule] (0,0) circle (1.17) (0,0) circle (1.0);
    \draw[tpc bath edge] (0,0) circle (1.17);
    \draw[tpc axis, -{Stealth[length=4pt, width=3pt]}] (-1.25,0) -- (1.36,0);
    \draw[tpc axis, -{Stealth[length=4pt, width=3pt]}] (0,-1.22) -- (0,1.34);
    \draw[tpc axis, line width=\tpcstrokeregular] (0,0) circle (1.0);
    \draw[tpc dynamics, -{Stealth[length=5pt]}] (58:0.75) arc (58:112:0.75);
    \node[tpc label, inner sep=1pt] at (136:0.56) {$U_t$};
    \draw[tpc effect, solid, line width=2.2pt] (0,0) -- (0.482,0);
    \draw[tpc effect, -{Stealth[length=5pt]}] (0.482,0) -- (1.0,0)
      node[tpc label, below right, inner sep=1.5pt] {$\vn$};
    \draw[tpc leader, dash pattern=on 1.6pt off 1.4pt] (50:0.75) -- (0.482,0);
    \draw[tpc state, -{Stealth[length=5pt]}] (0,0) -- (50:0.75)
      node[tpc label, right, inner sep=2.5pt] {$\vr$};
    \node[tpc label, below, inner sep=3pt] at (0.24,0) {$\vr\cdot\vn$};
    \node[tpc label, anchor=north west, inner sep=0pt] at (-1.24,1.36) {(a)};
    \node[tpc label, anchor=north west, align=left] at (0.72,-0.86) {keine\\[-2pt]Zustände};
  \end{scope}
  % (b) ladder: E_n = n^2 E_1, n = 1..6; 0.075 cm per E_1
  \begin{scope}[xshift=3.2cm, yshift=-1.5cm]
    \node[tpc label, anchor=north west, inner sep=0pt] at (-0.55,3.24) {(b)};
    \draw[tpc axis, -{Stealth[length=4pt, width=3pt]}] (0.2,-0.1) -- (0.2,3.05)
      node[tpc label, right, inner sep=2pt] {$E$};
    \foreach \n in {1,...,6}{
      \draw[tpc q\n] (0.45,{0.075*\n*\n}) -- (1.75,{0.075*\n*\n});
      \node[tpc label, right, inner sep=2pt] at (1.78,{0.075*\n*\n+(\n==1 ? -0.07 : 0)+(\n==2 ? 0.03 : 0)}) {$\n$};}
    \draw[tpc dynamics, -{Stealth[length=5pt]}] (1.1,{0.075*9}) -- (1.1,{0.075*4+0.03});
    \node[tpc label, anchor=north west, inner sep=1pt] at (0.42,-0.12) {$E_n=n^2E_1$};
  \end{scope}
\end{tikzpicture}}%
\hspace{\tpcspaces}%
% (c) the same roles on the stage: what a 3D demo or a video frame draws with
{\fboxsep=\tpcplatepad\colorbox{tpc-stage}{%
\begin{tikzpicture}[x=1cm, y=1cm, baseline=(current bounding box.south)]
  \useasboundingbox (-0.1,-1.78) rectangle (3.5,1.78);
  \node[tpc stage label, anchor=north west, inner sep=0pt] at (-0.02,1.74) {(c) Bühne};
  \fill[tpc stage bath] (0,-0.2) rectangle (1.15,0.05);
  \fill[tpc stage bath hatch] (0,-0.2) rectangle (1.15,0.05);
  \draw[tpc stage bath edge] (0,-0.2) rectangle (1.15,0.05);
  \draw[tpc stage dynamics] (0,0.35) -- (1.15,0.35);
  \draw[tpc stage effect] (0,0.7) -- (1.15,0.7);
  \draw[tpc stage state] (0,1.05) -- (1.15,1.05);
  \foreach \t/\y in {Zustand/1.05, Effekt/0.7, Dynamik/0.35, Bad/-0.08}
    \node[tpc stage label, right, inner sep=3pt] at (1.15,\y) {\t};
  \foreach \n in {1,...,6}
    \draw[tpc stage q\n, line width=3.2pt] ({0.5*(\n-1)},-0.95) -- ({0.5*(\n-1)+0.42},-0.95);
  \node[tpc stage label, anchor=north west, inner sep=0pt] at (0,-1.15) {Leiter $n=1,\dots,6$};
\end{tikzpicture}}}%
}\par\vspace{\tpcspacexs}
\noindent{\tpcfontsmall\color{tpc-ink-muted}{\tpcfontlabel\bfseries\color{tpc-ink} Abb.~6.1}\enspace
\textit{Die Wahrscheinlichkeit ist eine Projektion.} (a) Zustand $\vr$, Detektorrichtung $\vn$,
Projektion $p=\tfrac12(1+\vr\cdot\vn)\approx 0{,}74$; außerhalb des Kreises liegen keine Zustände.
(b) Teilchen im Kasten, $n=1,\dots,6$. (c) Dieselben Rollen auf der Bühne für 3D und Video.\par}

\vspace{\tpcspacexs}
\begin{specvertiefung}{6.V2}{Zustände als Funktionale}{Fundament 6.F1: Spurklasse und
normale Zustände}
\noindent Voraussetzung: Kern 6.1. Allgemein ist ein Zustand ein positives, normiertes lineares
Funktional $\omega$ auf der Algebra $\mathcal{B}(\mathcal{H})$, also $\omega(A^{*}A)\geq 0$ und
$\omega(\one)=1$. Die \emph{normalen} Zustände sind genau die Funktionale
$\omega(A)=\operatorname{tr}(\rho A)$ mit einem Dichteoperator $\rho$ der Spurklasse.
{\color{tpc-depth-v}Zurück zum Kern:} im Folgenden genügen Dichteoperatoren.
\end{specvertiefung}

\specsection{Was diese Seite zeigt}
Schrift, Farbe und Maße dieser Seite stammen aus einer einzigen Tokendatei. Farbe trägt nur
Physik: Zustand, Effekt, Dynamik, Bad und die Leiter der Quantenzahlen.

\vfill
\noindent{\tpcfontmono\color{tpc-ink-muted}\tpctoktitle\\
pkg=\tpcdesignpkg\ ground=\tpcdesignground\ math=\tpcdesignmath\ fonts=\tpcdesignfonts}
\end{document}
"
% Every colour, face and length below is a token name from design/README.md.
% This file contains no package-specific value; design/lint_content.py checks that.
%
% Text and physics are taken from the lane specimen
% docs/plan/research/design-packages-A/specimens/body-de.tex (checked there by hand):
%   rho = (1 + r.sigma)/2, E = (1 + n.sigma)/2, |n| = 1  =>  tr(rho E) = (1 + r.n)/2.
% Figure geometry: r = 0.75 at 50 degrees, n = e_x, so r.n = 0.75 cos 50 = 0.482 and
% p = (1 + 0.482)/2 = 0.741.  Ladder: particle in a box, E_n = n^2 E_1
% (latex/lectures/de/vl11.tex, latex/figures/fig_vl11_kasten.tex), n = 1..6.
\providecommand\tpcspecpkg{1}
\providecommand\tpcspecground{light}
\providecommand\tpcspecmath{stix}
\providecommand\tpcspecfonts{open}
\documentclass[11pt]{article}
\usepackage[ngerman]{babel}
\usepackage[autostyle]{csquotes}
\usepackage{tikz}
\usetikzlibrary{arrows.meta}
\usepackage[pkg=\tpcspecpkg, ground=\tpcspecground, math=\tpcspecmath,
            fonts=\tpcspecfonts]{tpc-design}
\usepackage[most]{tcolorbox}
\usepackage[protrusion=true, expansion=false]{microtype}
\usepackage[hidelinks]{hyperref}
\tpcapplypage
\pagestyle{empty}
\emergencystretch=1em
\setlength{\parindent}{\tpcspacem}
\setlength{\parskip}{0pt}

% ---- semantic devices, built from tokens only ------------------------------
\newcommand\speckicker[1]{{\tpcfontlabel\color{tpc-ink-muted}#1}\par\vspace{\tpcspacexs}}
\newcommand\spechead[2]{%
  \noindent\llap{{\tpcfontnumeral\color{tpc-accent}#1}\hspace{\tpcspaces}}%
  {\tpcfonthi\color{tpc-accent}#2}\par\vspace{\tpcspaces}}
\newcommand\specsection[1]{\par\vspace{\tpcspaces}\noindent
  {\tpcfonthii\color{tpc-accent}#1}\par\vspace{\tpcspacexs}\noindent\ignorespaces}
\newcommand\specnote[1]{\marginpar{\raggedright\tpcfontlabel\color{tpc-ink-muted}#1}}
\newcommand\speclink[1]{{\color{tpc-link}\underline{#1}}}
% Merksatz: highlight tint, box rule in the accent, run-in label.
\newtcolorbox{specmerksatz}[1]{enhanced, frame hidden, boxrule=0pt, arc=0pt,
  colback=tpc-highlight, coltext=tpc-ink, left=\tpcspaces, right=\tpcspaces,
  top=\tpcspaces, bottom=\tpcspaces, before skip=\tpcspaces, after skip=\tpcspaces,
  borderline west={\tpcstrokebox}{0pt}{tpc-accent},
  boxsep=0pt,
  before upper={\begingroup\tpcfonthiii\color{tpc-accent}#1\endgroup\quad}}
% Depth marking, the same device in every package: body size, no tint, no small
% type; a rule in the outer margin over the block's length and a run-in label.
% Only the rule's colour and weight and the label's face come from the package.
% #1 label, #2 title, #3 pointer to the Fundament (set beside the block).
\newtcolorbox{specvertiefung}[3]{enhanced, frame hidden, boxrule=0pt, arc=0pt,
  colback=tpc-ground, coltext=tpc-ink, boxsep=0pt, left=0pt, right=0pt, top=0pt, bottom=0pt,
  before skip=\tpcspaces, after skip=\tpcspaces,
  overlay={%
    \fill[tpc-depth-v] ([xshift=\tpcspaces]frame.north east)
      rectangle ([xshift=\dimexpr\tpcspaces+\tpcstrokedepthv\relax]frame.south east);
    % Fundament pointer: double hairline in the depth-f role, then the label
    \coordinate (F) at ([xshift=\tpcmargingap, yshift=-3pt]frame.north east);
    \draw[tpc-depth-f, line width=\tpcstrokedepthf] (F) -- ++(\tpcmarginwidth,0)
      ([yshift=-2pt]F) -- ++(\tpcmarginwidth,0);
    \node[anchor=north west, inner sep=0pt, text width=\tpcmarginwidth, align=flush left,
      text=tpc-depth-f, font=\tpcfontlabel] at ([yshift=-5pt]F) {#3};},
  before upper={\begingroup\tpcfontdepth\color{tpc-depth-v}Vertiefung #1\enspace #2\par\endgroup}}
\newcommand\vr{\symbfit{r}}
\newcommand\vn{\symbfit{n}}
\newcommand\vsigma{\symbfit{\sigma}}
\ifdefined\symbfit\else
  \renewcommand\vr{\boldsymbol{r}}\renewcommand\vn{\boldsymbol{n}}
  \renewcommand\vsigma{\boldsymbol{\sigma}}
\fi
\newcommand\one{\mathbb{1}}

\renewcommand\theequation{6.\arabic{equation}}
\begin{document}
\speckicker{Zoomstufe 1 \textendash\ ein einzelnes Quantensystem}
\spechead{6}{Die Bornsche Regel}

\noindent Eine physikalische Theorie liefert Daten für drei Dinge: Präparationen, Dynamik und
Messungen. Eine Präparation nennen wir \enquote{Zustand}; im Quantenfall beschreiben wir sie
durch einen Dichteoperator $\rho\geq 0$ mit $\operatorname{tr}\rho=1$.\specnote{Für ein Qubit ist
jeder Dichteoperator von der Form $\rho=\tfrac12(\one+\vr\cdot\vsigma)$ mit
$\lvert\vr\rvert\leq 1$. Reine Zustände liegen auf dem Rand.}
Ein Detektor, der \enquote{klick} macht oder schweigt, wird durch einen Effekt $E$ mit
$0\leq E\leq\one$ beschrieben. Die Theorie sagt nur Wahrscheinlichkeiten voraus, und zwar genau
eine Zahl für jedes Paar aus Zustand und Effekt (siehe \speclink{Definition D-6.1}).

\begin{specmerksatz}{Postulat P4}
Wird ein System im Zustand $\rho$ präpariert und mit dem Effekt $E$ getestet, so spricht der
Detektor mit der Wahrscheinlichkeit $p(E\,\vert\,\rho)=\operatorname{tr}(\rho E)$ an.
\end{specmerksatz}

\noindent Für ein Qubit lässt sich die Regel zeichnen (Abb.~6.1a). Mit
$\rho=\tfrac12(\one+\vr\cdot\vsigma)$ und dem Projektor $E=\tfrac12(\one+\vn\cdot\vsigma)$,
$\lvert\vn\rvert=1$, folgt
\begin{equation}
  \operatorname{tr}(\rho E)=\tfrac12\bigl(1+\vr\cdot\vn\bigr),
  \qquad
  \rho\;\mapsto\;U_t\,\rho\,U_t^{*},\quad U_t=\mathrm{e}^{-\mathrm{i}Ht/\hbar}.
\end{equation}
Die Wahrscheinlichkeit ist also die Projektion des Bloch-Vektors auf die Richtung des Detektors,
verschoben und halbiert. Die Zeitentwicklung dreht den Vektor; sie ändert $\lvert\vr\rvert$ nicht.

% ---- figure: plate (a, b) and stage (c) -------------------------------------
\par\vspace{\tpcspaces}
\noindent\makebox[\linewidth][l]{%
\tpcplate{%
\begin{tikzpicture}[x=1cm, y=1cm, baseline=(current bounding box.south)]
  \useasboundingbox (-1.62,-1.78) rectangle (5.3,1.78);
  % (a) Bloch disc: state, effect, dynamics, bath
  \begin{scope}[scale=1.28]
    \fill[tpc bath, even odd rule] (0,0) circle (1.17) (0,0) circle (1.0);
    \fill[tpc bath hatch, even odd rule] (0,0) circle (1.17) (0,0) circle (1.0);
    \draw[tpc bath edge] (0,0) circle (1.17);
    \draw[tpc axis, -{Stealth[length=4pt, width=3pt]}] (-1.25,0) -- (1.36,0);
    \draw[tpc axis, -{Stealth[length=4pt, width=3pt]}] (0,-1.22) -- (0,1.34);
    \draw[tpc axis, line width=\tpcstrokeregular] (0,0) circle (1.0);
    \draw[tpc dynamics, -{Stealth[length=5pt]}] (58:0.75) arc (58:112:0.75);
    \node[tpc label, inner sep=1pt] at (136:0.56) {$U_t$};
    \draw[tpc effect, solid, line width=2.2pt] (0,0) -- (0.482,0);
    \draw[tpc effect, -{Stealth[length=5pt]}] (0.482,0) -- (1.0,0)
      node[tpc label, below right, inner sep=1.5pt] {$\vn$};
    \draw[tpc leader, dash pattern=on 1.6pt off 1.4pt] (50:0.75) -- (0.482,0);
    \draw[tpc state, -{Stealth[length=5pt]}] (0,0) -- (50:0.75)
      node[tpc label, right, inner sep=2.5pt] {$\vr$};
    \node[tpc label, below, inner sep=3pt] at (0.24,0) {$\vr\cdot\vn$};
    \node[tpc label, anchor=north west, inner sep=0pt] at (-1.24,1.36) {(a)};
    \node[tpc label, anchor=north west, align=left] at (0.72,-0.86) {keine\\[-2pt]Zustände};
  \end{scope}
  % (b) ladder: E_n = n^2 E_1, n = 1..6; 0.075 cm per E_1
  \begin{scope}[xshift=3.2cm, yshift=-1.5cm]
    \node[tpc label, anchor=north west, inner sep=0pt] at (-0.55,3.24) {(b)};
    \draw[tpc axis, -{Stealth[length=4pt, width=3pt]}] (0.2,-0.1) -- (0.2,3.05)
      node[tpc label, right, inner sep=2pt] {$E$};
    \foreach \n in {1,...,6}{
      \draw[tpc q\n] (0.45,{0.075*\n*\n}) -- (1.75,{0.075*\n*\n});
      \node[tpc label, right, inner sep=2pt] at (1.78,{0.075*\n*\n+(\n==1 ? -0.07 : 0)+(\n==2 ? 0.03 : 0)}) {$\n$};}
    \draw[tpc dynamics, -{Stealth[length=5pt]}] (1.1,{0.075*9}) -- (1.1,{0.075*4+0.03});
    \node[tpc label, anchor=north west, inner sep=1pt] at (0.42,-0.12) {$E_n=n^2E_1$};
  \end{scope}
\end{tikzpicture}}%
\hspace{\tpcspaces}%
% (c) the same roles on the stage: what a 3D demo or a video frame draws with
{\fboxsep=\tpcplatepad\colorbox{tpc-stage}{%
\begin{tikzpicture}[x=1cm, y=1cm, baseline=(current bounding box.south)]
  \useasboundingbox (-0.1,-1.78) rectangle (3.5,1.78);
  \node[tpc stage label, anchor=north west, inner sep=0pt] at (-0.02,1.74) {(c) Bühne};
  \fill[tpc stage bath] (0,-0.2) rectangle (1.15,0.05);
  \fill[tpc stage bath hatch] (0,-0.2) rectangle (1.15,0.05);
  \draw[tpc stage bath edge] (0,-0.2) rectangle (1.15,0.05);
  \draw[tpc stage dynamics] (0,0.35) -- (1.15,0.35);
  \draw[tpc stage effect] (0,0.7) -- (1.15,0.7);
  \draw[tpc stage state] (0,1.05) -- (1.15,1.05);
  \foreach \t/\y in {Zustand/1.05, Effekt/0.7, Dynamik/0.35, Bad/-0.08}
    \node[tpc stage label, right, inner sep=3pt] at (1.15,\y) {\t};
  \foreach \n in {1,...,6}
    \draw[tpc stage q\n, line width=3.2pt] ({0.5*(\n-1)},-0.95) -- ({0.5*(\n-1)+0.42},-0.95);
  \node[tpc stage label, anchor=north west, inner sep=0pt] at (0,-1.15) {Leiter $n=1,\dots,6$};
\end{tikzpicture}}}%
}\par\vspace{\tpcspacexs}
\noindent{\tpcfontsmall\color{tpc-ink-muted}{\tpcfontlabel\bfseries\color{tpc-ink} Abb.~6.1}\enspace
\textit{Die Wahrscheinlichkeit ist eine Projektion.} (a) Zustand $\vr$, Detektorrichtung $\vn$,
Projektion $p=\tfrac12(1+\vr\cdot\vn)\approx 0{,}74$; außerhalb des Kreises liegen keine Zustände.
(b) Teilchen im Kasten, $n=1,\dots,6$. (c) Dieselben Rollen auf der Bühne für 3D und Video.\par}

\vspace{\tpcspacexs}
\begin{specvertiefung}{6.V2}{Zustände als Funktionale}{Fundament 6.F1: Spurklasse und
normale Zustände}
\noindent Voraussetzung: Kern 6.1. Allgemein ist ein Zustand ein positives, normiertes lineares
Funktional $\omega$ auf der Algebra $\mathcal{B}(\mathcal{H})$, also $\omega(A^{*}A)\geq 0$ und
$\omega(\one)=1$. Die \emph{normalen} Zustände sind genau die Funktionale
$\omega(A)=\operatorname{tr}(\rho A)$ mit einem Dichteoperator $\rho$ der Spurklasse.
{\color{tpc-depth-v}Zurück zum Kern:} im Folgenden genügen Dichteoperatoren.
\end{specvertiefung}

\specsection{Was diese Seite zeigt}
Schrift, Farbe und Maße dieser Seite stammen aus einer einzigen Tokendatei. Farbe trägt nur
Physik: Zustand, Effekt, Dynamik, Bad und die Leiter der Quantenzahlen.

\vfill
\noindent{\tpcfontmono\color{tpc-ink-muted}\tpctoktitle\\
pkg=\tpcdesignpkg\ ground=\tpcdesignground\ math=\tpcdesignmath\ fonts=\tpcdesignfonts}
\end{document}
"
% Every colour, face and length below is a token name from design/README.md.
% This file contains no package-specific value; design/lint_content.py checks that.
%
% Text and physics are taken from the lane specimen
% docs/plan/research/design-packages-A/specimens/body-de.tex (checked there by hand):
%   rho = (1 + r.sigma)/2, E = (1 + n.sigma)/2, |n| = 1  =>  tr(rho E) = (1 + r.n)/2.
% Figure geometry: r = 0.75 at 50 degrees, n = e_x, so r.n = 0.75 cos 50 = 0.482 and
% p = (1 + 0.482)/2 = 0.741.  Ladder: particle in a box, E_n = n^2 E_1
% (latex/lectures/de/vl11.tex, latex/figures/fig_vl11_kasten.tex), n = 1..6.
\providecommand\tpcspecpkg{1}
\providecommand\tpcspecground{light}
\providecommand\tpcspecmath{stix}
\providecommand\tpcspecfonts{open}
\documentclass[11pt]{article}
\usepackage[ngerman]{babel}
\usepackage[autostyle]{csquotes}
\usepackage{tikz}
\usetikzlibrary{arrows.meta}
\usepackage[pkg=\tpcspecpkg, ground=\tpcspecground, math=\tpcspecmath,
            fonts=\tpcspecfonts]{tpc-design}
\usepackage[most]{tcolorbox}
\usepackage[protrusion=true, expansion=false]{microtype}
\usepackage[hidelinks]{hyperref}
\tpcapplypage
\pagestyle{empty}
\emergencystretch=1em
\setlength{\parindent}{\tpcspacem}
\setlength{\parskip}{0pt}

% ---- semantic devices, built from tokens only ------------------------------
\newcommand\speckicker[1]{{\tpcfontlabel\color{tpc-ink-muted}#1}\par\vspace{\tpcspacexs}}
\newcommand\spechead[2]{%
  \noindent\llap{{\tpcfontnumeral\color{tpc-accent}#1}\hspace{\tpcspaces}}%
  {\tpcfonthi\color{tpc-accent}#2}\par\vspace{\tpcspaces}}
\newcommand\specsection[1]{\par\vspace{\tpcspaces}\noindent
  {\tpcfonthii\color{tpc-accent}#1}\par\vspace{\tpcspacexs}\noindent\ignorespaces}
\newcommand\specnote[1]{\marginpar{\raggedright\tpcfontlabel\color{tpc-ink-muted}#1}}
\newcommand\speclink[1]{{\color{tpc-link}\underline{#1}}}
% Merksatz: highlight tint, box rule in the accent, run-in label.
\newtcolorbox{specmerksatz}[1]{enhanced, frame hidden, boxrule=0pt, arc=0pt,
  colback=tpc-highlight, coltext=tpc-ink, left=\tpcspaces, right=\tpcspaces,
  top=\tpcspaces, bottom=\tpcspaces, before skip=\tpcspaces, after skip=\tpcspaces,
  borderline west={\tpcstrokebox}{0pt}{tpc-accent},
  boxsep=0pt,
  before upper={\begingroup\tpcfonthiii\color{tpc-accent}#1\endgroup\quad}}
% Depth marking, the same device in every package: body size, no tint, no small
% type; a rule in the outer margin over the block's length and a run-in label.
% Only the rule's colour and weight and the label's face come from the package.
% #1 label, #2 title, #3 pointer to the Fundament (set beside the block).
\newtcolorbox{specvertiefung}[3]{enhanced, frame hidden, boxrule=0pt, arc=0pt,
  colback=tpc-ground, coltext=tpc-ink, boxsep=0pt, left=0pt, right=0pt, top=0pt, bottom=0pt,
  before skip=\tpcspaces, after skip=\tpcspaces,
  overlay={%
    \fill[tpc-depth-v] ([xshift=\tpcspaces]frame.north east)
      rectangle ([xshift=\dimexpr\tpcspaces+\tpcstrokedepthv\relax]frame.south east);
    % Fundament pointer: double hairline in the depth-f role, then the label
    \coordinate (F) at ([xshift=\tpcmargingap, yshift=-3pt]frame.north east);
    \draw[tpc-depth-f, line width=\tpcstrokedepthf] (F) -- ++(\tpcmarginwidth,0)
      ([yshift=-2pt]F) -- ++(\tpcmarginwidth,0);
    \node[anchor=north west, inner sep=0pt, text width=\tpcmarginwidth, align=flush left,
      text=tpc-depth-f, font=\tpcfontlabel] at ([yshift=-5pt]F) {#3};},
  before upper={\begingroup\tpcfontdepth\color{tpc-depth-v}Vertiefung #1\enspace #2\par\endgroup}}
\newcommand\vr{\symbfit{r}}
\newcommand\vn{\symbfit{n}}
\newcommand\vsigma{\symbfit{\sigma}}
\ifdefined\symbfit\else
  \renewcommand\vr{\boldsymbol{r}}\renewcommand\vn{\boldsymbol{n}}
  \renewcommand\vsigma{\boldsymbol{\sigma}}
\fi
\newcommand\one{\mathbb{1}}

\renewcommand\theequation{6.\arabic{equation}}
\begin{document}
\speckicker{Zoomstufe 1 \textendash\ ein einzelnes Quantensystem}
\spechead{6}{Die Bornsche Regel}

\noindent Eine physikalische Theorie liefert Daten für drei Dinge: Präparationen, Dynamik und
Messungen. Eine Präparation nennen wir \enquote{Zustand}; im Quantenfall beschreiben wir sie
durch einen Dichteoperator $\rho\geq 0$ mit $\operatorname{tr}\rho=1$.\specnote{Für ein Qubit ist
jeder Dichteoperator von der Form $\rho=\tfrac12(\one+\vr\cdot\vsigma)$ mit
$\lvert\vr\rvert\leq 1$. Reine Zustände liegen auf dem Rand.}
Ein Detektor, der \enquote{klick} macht oder schweigt, wird durch einen Effekt $E$ mit
$0\leq E\leq\one$ beschrieben. Die Theorie sagt nur Wahrscheinlichkeiten voraus, und zwar genau
eine Zahl für jedes Paar aus Zustand und Effekt (siehe \speclink{Definition D-6.1}).

\begin{specmerksatz}{Postulat P4}
Wird ein System im Zustand $\rho$ präpariert und mit dem Effekt $E$ getestet, so spricht der
Detektor mit der Wahrscheinlichkeit $p(E\,\vert\,\rho)=\operatorname{tr}(\rho E)$ an.
\end{specmerksatz}

\noindent Für ein Qubit lässt sich die Regel zeichnen (Abb.~6.1a). Mit
$\rho=\tfrac12(\one+\vr\cdot\vsigma)$ und dem Projektor $E=\tfrac12(\one+\vn\cdot\vsigma)$,
$\lvert\vn\rvert=1$, folgt
\begin{equation}
  \operatorname{tr}(\rho E)=\tfrac12\bigl(1+\vr\cdot\vn\bigr),
  \qquad
  \rho\;\mapsto\;U_t\,\rho\,U_t^{*},\quad U_t=\mathrm{e}^{-\mathrm{i}Ht/\hbar}.
\end{equation}
Die Wahrscheinlichkeit ist also die Projektion des Bloch-Vektors auf die Richtung des Detektors,
verschoben und halbiert. Die Zeitentwicklung dreht den Vektor; sie ändert $\lvert\vr\rvert$ nicht.

% ---- figure: plate (a, b) and stage (c) -------------------------------------
\par\vspace{\tpcspaces}
\noindent\makebox[\linewidth][l]{%
\tpcplate{%
\begin{tikzpicture}[x=1cm, y=1cm, baseline=(current bounding box.south)]
  \useasboundingbox (-1.62,-1.78) rectangle (5.3,1.78);
  % (a) Bloch disc: state, effect, dynamics, bath
  \begin{scope}[scale=1.28]
    \fill[tpc bath, even odd rule] (0,0) circle (1.17) (0,0) circle (1.0);
    \fill[tpc bath hatch, even odd rule] (0,0) circle (1.17) (0,0) circle (1.0);
    \draw[tpc bath edge] (0,0) circle (1.17);
    \draw[tpc axis, -{Stealth[length=4pt, width=3pt]}] (-1.25,0) -- (1.36,0);
    \draw[tpc axis, -{Stealth[length=4pt, width=3pt]}] (0,-1.22) -- (0,1.34);
    \draw[tpc axis, line width=\tpcstrokeregular] (0,0) circle (1.0);
    \draw[tpc dynamics, -{Stealth[length=5pt]}] (58:0.75) arc (58:112:0.75);
    \node[tpc label, inner sep=1pt] at (136:0.56) {$U_t$};
    \draw[tpc effect, solid, line width=2.2pt] (0,0) -- (0.482,0);
    \draw[tpc effect, -{Stealth[length=5pt]}] (0.482,0) -- (1.0,0)
      node[tpc label, below right, inner sep=1.5pt] {$\vn$};
    \draw[tpc leader, dash pattern=on 1.6pt off 1.4pt] (50:0.75) -- (0.482,0);
    \draw[tpc state, -{Stealth[length=5pt]}] (0,0) -- (50:0.75)
      node[tpc label, right, inner sep=2.5pt] {$\vr$};
    \node[tpc label, below, inner sep=3pt] at (0.24,0) {$\vr\cdot\vn$};
    \node[tpc label, anchor=north west, inner sep=0pt] at (-1.24,1.36) {(a)};
    \node[tpc label, anchor=north west, align=left] at (0.72,-0.86) {keine\\[-2pt]Zustände};
  \end{scope}
  % (b) ladder: E_n = n^2 E_1, n = 1..6; 0.075 cm per E_1
  \begin{scope}[xshift=3.2cm, yshift=-1.5cm]
    \node[tpc label, anchor=north west, inner sep=0pt] at (-0.55,3.24) {(b)};
    \draw[tpc axis, -{Stealth[length=4pt, width=3pt]}] (0.2,-0.1) -- (0.2,3.05)
      node[tpc label, right, inner sep=2pt] {$E$};
    \foreach \n in {1,...,6}{
      \draw[tpc q\n] (0.45,{0.075*\n*\n}) -- (1.75,{0.075*\n*\n});
      \node[tpc label, right, inner sep=2pt] at (1.78,{0.075*\n*\n+(\n==1 ? -0.07 : 0)+(\n==2 ? 0.03 : 0)}) {$\n$};}
    \draw[tpc dynamics, -{Stealth[length=5pt]}] (1.1,{0.075*9}) -- (1.1,{0.075*4+0.03});
    \node[tpc label, anchor=north west, inner sep=1pt] at (0.42,-0.12) {$E_n=n^2E_1$};
  \end{scope}
\end{tikzpicture}}%
\hspace{\tpcspaces}%
% (c) the same roles on the stage: what a 3D demo or a video frame draws with
{\fboxsep=\tpcplatepad\colorbox{tpc-stage}{%
\begin{tikzpicture}[x=1cm, y=1cm, baseline=(current bounding box.south)]
  \useasboundingbox (-0.1,-1.78) rectangle (3.5,1.78);
  \node[tpc stage label, anchor=north west, inner sep=0pt] at (-0.02,1.74) {(c) Bühne};
  \fill[tpc stage bath] (0,-0.2) rectangle (1.15,0.05);
  \fill[tpc stage bath hatch] (0,-0.2) rectangle (1.15,0.05);
  \draw[tpc stage bath edge] (0,-0.2) rectangle (1.15,0.05);
  \draw[tpc stage dynamics] (0,0.35) -- (1.15,0.35);
  \draw[tpc stage effect] (0,0.7) -- (1.15,0.7);
  \draw[tpc stage state] (0,1.05) -- (1.15,1.05);
  \foreach \t/\y in {Zustand/1.05, Effekt/0.7, Dynamik/0.35, Bad/-0.08}
    \node[tpc stage label, right, inner sep=3pt] at (1.15,\y) {\t};
  \foreach \n in {1,...,6}
    \draw[tpc stage q\n, line width=3.2pt] ({0.5*(\n-1)},-0.95) -- ({0.5*(\n-1)+0.42},-0.95);
  \node[tpc stage label, anchor=north west, inner sep=0pt] at (0,-1.15) {Leiter $n=1,\dots,6$};
\end{tikzpicture}}}%
}\par\vspace{\tpcspacexs}
\noindent{\tpcfontsmall\color{tpc-ink-muted}{\tpcfontlabel\bfseries\color{tpc-ink} Abb.~6.1}\enspace
\textit{Die Wahrscheinlichkeit ist eine Projektion.} (a) Zustand $\vr$, Detektorrichtung $\vn$,
Projektion $p=\tfrac12(1+\vr\cdot\vn)\approx 0{,}74$; außerhalb des Kreises liegen keine Zustände.
(b) Teilchen im Kasten, $n=1,\dots,6$. (c) Dieselben Rollen auf der Bühne für 3D und Video.\par}

\vspace{\tpcspacexs}
\begin{specvertiefung}{6.V2}{Zustände als Funktionale}{Fundament 6.F1: Spurklasse und
normale Zustände}
\noindent Voraussetzung: Kern 6.1. Allgemein ist ein Zustand ein positives, normiertes lineares
Funktional $\omega$ auf der Algebra $\mathcal{B}(\mathcal{H})$, also $\omega(A^{*}A)\geq 0$ und
$\omega(\one)=1$. Die \emph{normalen} Zustände sind genau die Funktionale
$\omega(A)=\operatorname{tr}(\rho A)$ mit einem Dichteoperator $\rho$ der Spurklasse.
{\color{tpc-depth-v}Zurück zum Kern:} im Folgenden genügen Dichteoperatoren.
\end{specvertiefung}

\specsection{Was diese Seite zeigt}
Schrift, Farbe und Maße dieser Seite stammen aus einer einzigen Tokendatei. Farbe trägt nur
Physik: Zustand, Effekt, Dynamik, Bad und die Leiter der Quantenzahlen.

\vfill
\noindent{\tpcfontmono\color{tpc-ink-muted}\tpctoktitle\\
pkg=\tpcdesignpkg\ ground=\tpcdesignground\ math=\tpcdesignmath\ fonts=\tpcdesignfonts}
\end{document}
"
% Every colour, face and length below is a token name from design/README.md.
% This file contains no package-specific value; design/lint_content.py checks that.
%
% Text and physics are taken from the lane specimen
% docs/plan/research/design-packages-A/specimens/body-de.tex (checked there by hand):
%   rho = (1 + r.sigma)/2, E = (1 + n.sigma)/2, |n| = 1  =>  tr(rho E) = (1 + r.n)/2.
% Figure geometry: r = 0.75 at 50 degrees, n = e_x, so r.n = 0.75 cos 50 = 0.482 and
% p = (1 + 0.482)/2 = 0.741.  Ladder: particle in a box, E_n = n^2 E_1
% (latex/lectures/de/vl11.tex, latex/figures/fig_vl11_kasten.tex), n = 1..6.
\providecommand\tpcspecpkg{1}
\providecommand\tpcspecground{light}
\providecommand\tpcspecmath{stix}
\providecommand\tpcspecfonts{open}
\documentclass[11pt]{article}
\usepackage[ngerman]{babel}
\usepackage[autostyle]{csquotes}
\usepackage{tikz}
\usetikzlibrary{arrows.meta}
\usepackage[pkg=\tpcspecpkg, ground=\tpcspecground, math=\tpcspecmath,
            fonts=\tpcspecfonts]{tpc-design}
\usepackage[most]{tcolorbox}
\usepackage[protrusion=true, expansion=false]{microtype}
\usepackage[hidelinks]{hyperref}
\tpcapplypage
\pagestyle{empty}
\emergencystretch=1em
\setlength{\parindent}{\tpcspacem}
\setlength{\parskip}{0pt}

% ---- semantic devices, built from tokens only ------------------------------
\newcommand\speckicker[1]{{\tpcfontlabel\color{tpc-ink-muted}#1}\par\vspace{\tpcspacexs}}
\newcommand\spechead[2]{%
  \noindent\llap{{\tpcfontnumeral\color{tpc-accent}#1}\hspace{\tpcspaces}}%
  {\tpcfonthi\color{tpc-accent}#2}\par\vspace{\tpcspaces}}
\newcommand\specsection[1]{\par\vspace{\tpcspaces}\noindent
  {\tpcfonthii\color{tpc-accent}#1}\par\vspace{\tpcspacexs}\noindent\ignorespaces}
\newcommand\specnote[1]{\marginpar{\raggedright\tpcfontlabel\color{tpc-ink-muted}#1}}
\newcommand\speclink[1]{{\color{tpc-link}\underline{#1}}}
% Merksatz: highlight tint, box rule in the accent, run-in label.
\newtcolorbox{specmerksatz}[1]{enhanced, frame hidden, boxrule=0pt, arc=0pt,
  colback=tpc-highlight, coltext=tpc-ink, left=\tpcspaces, right=\tpcspaces,
  top=\tpcspaces, bottom=\tpcspaces, before skip=\tpcspaces, after skip=\tpcspaces,
  borderline west={\tpcstrokebox}{0pt}{tpc-accent},
  boxsep=0pt,
  before upper={\begingroup\tpcfonthiii\color{tpc-accent}#1\endgroup\quad}}
% Depth marking, the same device in every package: body size, no tint, no small
% type; a rule in the outer margin over the block's length and a run-in label.
% Only the rule's colour and weight and the label's face come from the package.
% #1 label, #2 title, #3 pointer to the Fundament (set beside the block).
\newtcolorbox{specvertiefung}[3]{enhanced, frame hidden, boxrule=0pt, arc=0pt,
  colback=tpc-ground, coltext=tpc-ink, boxsep=0pt, left=0pt, right=0pt, top=0pt, bottom=0pt,
  before skip=\tpcspaces, after skip=\tpcspaces,
  overlay={%
    \fill[tpc-depth-v] ([xshift=\tpcspaces]frame.north east)
      rectangle ([xshift=\dimexpr\tpcspaces+\tpcstrokedepthv\relax]frame.south east);
    % Fundament pointer: double hairline in the depth-f role, then the label
    \coordinate (F) at ([xshift=\tpcmargingap, yshift=-3pt]frame.north east);
    \draw[tpc-depth-f, line width=\tpcstrokedepthf] (F) -- ++(\tpcmarginwidth,0)
      ([yshift=-2pt]F) -- ++(\tpcmarginwidth,0);
    \node[anchor=north west, inner sep=0pt, text width=\tpcmarginwidth, align=flush left,
      text=tpc-depth-f, font=\tpcfontlabel] at ([yshift=-5pt]F) {#3};},
  before upper={\begingroup\tpcfontdepth\color{tpc-depth-v}Vertiefung #1\enspace #2\par\endgroup}}
\newcommand\vr{\symbfit{r}}
\newcommand\vn{\symbfit{n}}
\newcommand\vsigma{\symbfit{\sigma}}
\ifdefined\symbfit\else
  \renewcommand\vr{\boldsymbol{r}}\renewcommand\vn{\boldsymbol{n}}
  \renewcommand\vsigma{\boldsymbol{\sigma}}
\fi
\newcommand\one{\mathbb{1}}

\renewcommand\theequation{6.\arabic{equation}}
\begin{document}
\speckicker{Zoomstufe 1 \textendash\ ein einzelnes Quantensystem}
\spechead{6}{Die Bornsche Regel}

\noindent Eine physikalische Theorie liefert Daten für drei Dinge: Präparationen, Dynamik und
Messungen. Eine Präparation nennen wir \enquote{Zustand}; im Quantenfall beschreiben wir sie
durch einen Dichteoperator $\rho\geq 0$ mit $\operatorname{tr}\rho=1$.\specnote{Für ein Qubit ist
jeder Dichteoperator von der Form $\rho=\tfrac12(\one+\vr\cdot\vsigma)$ mit
$\lvert\vr\rvert\leq 1$. Reine Zustände liegen auf dem Rand.}
Ein Detektor, der \enquote{klick} macht oder schweigt, wird durch einen Effekt $E$ mit
$0\leq E\leq\one$ beschrieben. Die Theorie sagt nur Wahrscheinlichkeiten voraus, und zwar genau
eine Zahl für jedes Paar aus Zustand und Effekt (siehe \speclink{Definition D-6.1}).

\begin{specmerksatz}{Postulat P4}
Wird ein System im Zustand $\rho$ präpariert und mit dem Effekt $E$ getestet, so spricht der
Detektor mit der Wahrscheinlichkeit $p(E\,\vert\,\rho)=\operatorname{tr}(\rho E)$ an.
\end{specmerksatz}

\noindent Für ein Qubit lässt sich die Regel zeichnen (Abb.~6.1a). Mit
$\rho=\tfrac12(\one+\vr\cdot\vsigma)$ und dem Projektor $E=\tfrac12(\one+\vn\cdot\vsigma)$,
$\lvert\vn\rvert=1$, folgt
\begin{equation}
  \operatorname{tr}(\rho E)=\tfrac12\bigl(1+\vr\cdot\vn\bigr),
  \qquad
  \rho\;\mapsto\;U_t\,\rho\,U_t^{*},\quad U_t=\mathrm{e}^{-\mathrm{i}Ht/\hbar}.
\end{equation}
Die Wahrscheinlichkeit ist also die Projektion des Bloch-Vektors auf die Richtung des Detektors,
verschoben und halbiert. Die Zeitentwicklung dreht den Vektor; sie ändert $\lvert\vr\rvert$ nicht.

% ---- figure: plate (a, b) and stage (c) -------------------------------------
\par\vspace{\tpcspaces}
\noindent\makebox[\linewidth][l]{%
\tpcplate{%
\begin{tikzpicture}[x=1cm, y=1cm, baseline=(current bounding box.south)]
  \useasboundingbox (-1.62,-1.78) rectangle (5.3,1.78);
  % (a) Bloch disc: state, effect, dynamics, bath
  \begin{scope}[scale=1.28]
    \fill[tpc bath, even odd rule] (0,0) circle (1.17) (0,0) circle (1.0);
    \fill[tpc bath hatch, even odd rule] (0,0) circle (1.17) (0,0) circle (1.0);
    \draw[tpc bath edge] (0,0) circle (1.17);
    \draw[tpc axis, -{Stealth[length=4pt, width=3pt]}] (-1.25,0) -- (1.36,0);
    \draw[tpc axis, -{Stealth[length=4pt, width=3pt]}] (0,-1.22) -- (0,1.34);
    \draw[tpc axis, line width=\tpcstrokeregular] (0,0) circle (1.0);
    \draw[tpc dynamics, -{Stealth[length=5pt]}] (58:0.75) arc (58:112:0.75);
    \node[tpc label, inner sep=1pt] at (136:0.56) {$U_t$};
    \draw[tpc effect, solid, line width=2.2pt] (0,0) -- (0.482,0);
    \draw[tpc effect, -{Stealth[length=5pt]}] (0.482,0) -- (1.0,0)
      node[tpc label, below right, inner sep=1.5pt] {$\vn$};
    \draw[tpc leader, dash pattern=on 1.6pt off 1.4pt] (50:0.75) -- (0.482,0);
    \draw[tpc state, -{Stealth[length=5pt]}] (0,0) -- (50:0.75)
      node[tpc label, right, inner sep=2.5pt] {$\vr$};
    \node[tpc label, below, inner sep=3pt] at (0.24,0) {$\vr\cdot\vn$};
    \node[tpc label, anchor=north west, inner sep=0pt] at (-1.24,1.36) {(a)};
    \node[tpc label, anchor=north west, align=left] at (0.72,-0.86) {keine\\[-2pt]Zustände};
  \end{scope}
  % (b) ladder: E_n = n^2 E_1, n = 1..6; 0.075 cm per E_1
  \begin{scope}[xshift=3.2cm, yshift=-1.5cm]
    \node[tpc label, anchor=north west, inner sep=0pt] at (-0.55,3.24) {(b)};
    \draw[tpc axis, -{Stealth[length=4pt, width=3pt]}] (0.2,-0.1) -- (0.2,3.05)
      node[tpc label, right, inner sep=2pt] {$E$};
    \foreach \n in {1,...,6}{
      \draw[tpc q\n] (0.45,{0.075*\n*\n}) -- (1.75,{0.075*\n*\n});
      \node[tpc label, right, inner sep=2pt] at (1.78,{0.075*\n*\n+(\n==1 ? -0.07 : 0)+(\n==2 ? 0.03 : 0)}) {$\n$};}
    \draw[tpc dynamics, -{Stealth[length=5pt]}] (1.1,{0.075*9}) -- (1.1,{0.075*4+0.03});
    \node[tpc label, anchor=north west, inner sep=1pt] at (0.42,-0.12) {$E_n=n^2E_1$};
  \end{scope}
\end{tikzpicture}}%
\hspace{\tpcspaces}%
% (c) the same roles on the stage: what a 3D demo or a video frame draws with
{\fboxsep=\tpcplatepad\colorbox{tpc-stage}{%
\begin{tikzpicture}[x=1cm, y=1cm, baseline=(current bounding box.south)]
  \useasboundingbox (-0.1,-1.78) rectangle (3.5,1.78);
  \node[tpc stage label, anchor=north west, inner sep=0pt] at (-0.02,1.74) {(c) Bühne};
  \fill[tpc stage bath] (0,-0.2) rectangle (1.15,0.05);
  \fill[tpc stage bath hatch] (0,-0.2) rectangle (1.15,0.05);
  \draw[tpc stage bath edge] (0,-0.2) rectangle (1.15,0.05);
  \draw[tpc stage dynamics] (0,0.35) -- (1.15,0.35);
  \draw[tpc stage effect] (0,0.7) -- (1.15,0.7);
  \draw[tpc stage state] (0,1.05) -- (1.15,1.05);
  \foreach \t/\y in {Zustand/1.05, Effekt/0.7, Dynamik/0.35, Bad/-0.08}
    \node[tpc stage label, right, inner sep=3pt] at (1.15,\y) {\t};
  \foreach \n in {1,...,6}
    \draw[tpc stage q\n, line width=3.2pt] ({0.5*(\n-1)},-0.95) -- ({0.5*(\n-1)+0.42},-0.95);
  \node[tpc stage label, anchor=north west, inner sep=0pt] at (0,-1.15) {Leiter $n=1,\dots,6$};
\end{tikzpicture}}}%
}\par\vspace{\tpcspacexs}
\noindent{\tpcfontsmall\color{tpc-ink-muted}{\tpcfontlabel\bfseries\color{tpc-ink} Abb.~6.1}\enspace
\textit{Die Wahrscheinlichkeit ist eine Projektion.} (a) Zustand $\vr$, Detektorrichtung $\vn$,
Projektion $p=\tfrac12(1+\vr\cdot\vn)\approx 0{,}74$; außerhalb des Kreises liegen keine Zustände.
(b) Teilchen im Kasten, $n=1,\dots,6$. (c) Dieselben Rollen auf der Bühne für 3D und Video.\par}

\vspace{\tpcspacexs}
\begin{specvertiefung}{6.V2}{Zustände als Funktionale}{Fundament 6.F1: Spurklasse und
normale Zustände}
\noindent Voraussetzung: Kern 6.1. Allgemein ist ein Zustand ein positives, normiertes lineares
Funktional $\omega$ auf der Algebra $\mathcal{B}(\mathcal{H})$, also $\omega(A^{*}A)\geq 0$ und
$\omega(\one)=1$. Die \emph{normalen} Zustände sind genau die Funktionale
$\omega(A)=\operatorname{tr}(\rho A)$ mit einem Dichteoperator $\rho$ der Spurklasse.
{\color{tpc-depth-v}Zurück zum Kern:} im Folgenden genügen Dichteoperatoren.
\end{specvertiefung}

\specsection{Was diese Seite zeigt}
Schrift, Farbe und Maße dieser Seite stammen aus einer einzigen Tokendatei. Farbe trägt nur
Physik: Zustand, Effekt, Dynamik, Bad und die Leiter der Quantenzahlen.

\vfill
\noindent{\tpcfontmono\color{tpc-ink-muted}\tpctoktitle\\
pkg=\tpcdesignpkg\ ground=\tpcdesignground\ math=\tpcdesignmath\ fonts=\tpcdesignfonts}
\end{document}
